\section*{Chapter \ref{ch:g12:cells-and-electrolysis} --- Electrochemical Cells and Electrolysis}

\begin{solution}{exo:g12:cells-and-electrolysis:1}
The zinc electrode is the \omterm{def:g12:cells-and-electrolysis:anode}{anode}, \ce{Zn -> Zn^{2+} + 2e-}; the copper
electrode is the \omterm{def:g12:cells-and-electrolysis:anode}{cathode}, \ce{Cu^{2+} + 2e- -> Cu}.
\end{solution}

\begin{solution}{exo:g12:cells-and-electrolysis:2}
% ledger: const:F
$Q = 0.20 \times 1800 = \qty{360}{C}$; $n(e^-) = 360/96485 =
\qty{3.7e-3}{mol}$.
\end{solution}

\begin{solution}{exo:g12:cells-and-electrolysis:3}
It closes the circuit and keeps each solution electrically neutral, its
\omterm{def:g9:ions:ion}{ions} moving into the compartments. Without it the circuit is open: no
current flows, the voltmeter reads zero and a lamp stays dark.
\end{solution}

\begin{solution}{exo:g12:cells-and-electrolysis:4}
$2.5 \times 3600 = \qty{9000}{C}$.
\end{solution}

\begin{solution}{exo:g12:cells-and-electrolysis:5}
No: the reaction goes in the direction opposite to the spontaneous one.
The generator supplies the energy.
\end{solution}

\begin{solution}{exo:g12:cells-and-electrolysis:6}
\omterm{def:g12:cells-and-electrolysis:anode}{Anode} (copper, $-$): \ce{Cu -> Cu^{2+} + 2e-}. \omterm{def:g12:cells-and-electrolysis:anode}{Cathode} (silver, $+$):
\ce{Ag+ + e- -> Ag}. Overall: \ce{Cu + 2Ag+ -> Cu^{2+} + 2Ag}.
\end{solution}

\begin{solution}{exo:g12:cells-and-electrolysis:7}
Three years are $3 \times 365 \times 24 = \qty{26280}{h}$;
$0.010 \times 26280 = \qty{263}{mA.h}$.
\end{solution}

\begin{solution}{exo:g12:cells-and-electrolysis:8}
% ledger: const:F, aw:Ag
$Q = 0.50 \times 2400 = \qty{1200}{C}$; $n(e^-) = n(\ce{Ag}) =
1200/96485 = \qty{0.0124}{mol}$; $m = 0.0124 \times 107.9 =
\qty{1.34}{g}$.
\end{solution}

\begin{solution}{exo:g12:cells-and-electrolysis:9}
The current goes in the wire from the copper ($+$) to the zinc ($-$),
opposite to the \omterm{def:g9:inside-the-atom:nucleus}{electrons}. \ce{K+} \omterm{def:g9:ions:ion}{ions} move towards the copper
compartment, which loses \ce{Cu^{2+}}; \ce{NO3-} \omterm{def:g9:ions:ion}{ions} towards the zinc
compartment, which gains \ce{Zn^{2+}}.
\end{solution}

\begin{solution}{exo:g12:cells-and-electrolysis:10}
In both the \omterm{def:g12:cells-and-electrolysis:anode}{anode} is where the \omterm{def:g11:redox:oxidation}{oxidation} happens. In a cell the \omterm{def:g12:cells-and-electrolysis:anode}{anode}
is $-$, the reaction is spontaneous and chemical energy becomes
electrical energy; in an \omterm{def:g12:cells-and-electrolysis:electrolysis}{electrolysis} the \omterm{def:g12:cells-and-electrolysis:anode}{anode} is joined to the $+$ of
the generator, the reaction is forced and electrical energy becomes
chemical energy.
\end{solution}

\begin{solution}{exo:g12:cells-and-electrolysis:11}
\qty{12}{mL}: the equation \ce{2H2O -> 2H2 + O2} gives two \omterm{def:g7:atoms-and-molecules:molecule}{molecules} of
dihydrogen for one of dioxygen, and equal amounts of gas take equal
volumes.
\end{solution}

\begin{solution}{exo:g12:cells-and-electrolysis:12}
% ledger: aw:Zn, const:F
$n(\ce{Zn}) = 5.0/65.4 = \qty{0.076}{mol}$;
$n(\ce{Cu^{2+}}) = 0.100 \times 0.50 = \qty{0.050}{mol}$; they react
one to one, so the copper \omterm{def:g9:ions:ion}{ions} run out first. $n(e^-) = 2 \times 0.050
= \qty{0.100}{mol}$, $Q = \qty{9.65e3}{C} = \qty{2.68}{A.h}$. At
\qty{50}{mA}: $2.68/0.050 = \qty{54}{h}$.
\end{solution}

\begin{solution}{exo:g12:cells-and-electrolysis:13}
$Q = 2.0 \times 3600 = \qty{7200}{C}$; $E = 7200 \times 1.5 =
\qty{1.08e4}{J}$, that is $1.5 \times 2.0 = \qty{3.0}{W.h}$.
\end{solution}

\begin{solution}{exo:g12:cells-and-electrolysis:14}
% ledger: const:F
$Q = \qty{7200}{C}$, $n(e^-) = \qty{0.0746}{mol}$. Two \omterm{def:g9:inside-the-atom:nucleus}{electrons} per
\ce{H2}: $n(\ce{H2}) = \qty{0.0373}{mol}$, $V = \qty{0.90}{L}$. Four per
\ce{O2}: $n(\ce{O2}) = \qty{0.0187}{mol}$, $V = \qty{0.45}{L}$.
\end{solution}

\begin{solution}{exo:g12:cells-and-electrolysis:15}
The copper dissolved at the \omterm{def:g12:cells-and-electrolysis:anode}{anode} equals the copper deposited, the same
\omterm{def:g9:inside-the-atom:nucleus}{electrons} passing at both: $2.84/3.00 = \qty{94.7}{\percent}$ copper.
\end{solution}

\begin{solution}{pb:g12:cells-and-electrolysis:1}
% ledger: const:F, const:NA, aw:Li, dch:CH4
\textbf{1.} The \omterm{def:g12:cells-and-electrolysis:anode}{anode}: lithium is oxidised there.

\textbf{2.} The graphite electrode, the \omterm{def:g12:cells-and-electrolysis:anode}{anode}.

\textbf{3.} \omterm{def:g9:inside-the-atom:nucleus}{Electrons} go through the phone from the graphite electrode to
the oxide electrode; lithium \omterm{def:g9:ions:ion}{ions} cross the electrolyte in the same
direction, from graphite to oxide.

\textbf{4.} Its reaction can be driven backwards by a charger.

\textbf{5.} \qty{4.000}{A.h}, that is $4.000 \times 3600 =
\qty{14400}{C}$.

\textbf{6.} $14400/96485 = \qty{0.1492}{mol}$.

\textbf{7.} $0.1492 \times \num{6.022e23} = \num{8.99e22}$ \omterm{def:g9:inside-the-atom:nucleus}{electrons}.

\textbf{8.} $4.000/0.25 = \qty{16}{h}$.

\textbf{9.} $0.20 \times 14400 = \qty{2880}{C}$.

\textbf{10.} $E = 14400 \times 3.7 = \qty{5.3e4}{J}$, about \qty{53}{kJ}.

\textbf{11.} $53280/3600 = \qty{14.8}{W.h}$.

\textbf{12.} \omterm{def:g4:what-a-fire-needs:burning}{Burning} \qty{1}{g} of methane releases a little more
energy (\qty{55.7}{kJ}) than a full phone battery stores.

\textbf{13.} $1000/14.8 = 67.6$: 67 full charges.

\textbf{14.} \ce{Li+ + e- -> Li}: a \omterm{def:g11:redox:oxidation}{reduction}.

\textbf{15.} The charger forces the reaction in the direction opposite to
the spontaneous one.

\textbf{16.} $14400/2.0 = \qty{7200}{s}$, \qty{2.0}{h}.

\textbf{17.} One lithium \omterm{def:g9:ions:ion}{ion} per \omterm{def:g9:inside-the-atom:nucleus}{electron}: \qty{0.1492}{mol}.

\textbf{18.} $0.1492 \times 6.9 = \qty{1.03}{g}$.

\textbf{19.} About \qty{1.03}{g} of lithium shuttles between the
electrodes in one full charge.
\end{solution}
